
几何空间往往恰恰要借助其上的（良好的）函数才能得到最好的理解，这一点或许出人意料。例如，
$\mathbb{R}^n$ 的一个\trDifferentiable{}\trManifold{}便可以借其光滑（$C^\infty$）函数来研究。由于%
「几何空间」上处处有定义的函数可能寥寥无几，这一洞见的更精确版本是：空间的%
结构可以通过考察其全部\trOpenSubset{}上的全部函数而得到很好的把握。这些信息被编码并整理在
叫做\emph{\trSheaf}的东西之中。\trSheaf{}由\trLeray{}于二十世纪四十年代引入，将它引进\trAlgebraicGeometry{}的
则是\trSerre。（这个名字的由来将在\cref{thm:sheaf-name-origin}中略作解释。）我们将定义\trSheaf{}并陈述关于\trSheaf{}的
一些有用的事实。开篇先用一个动机性的例子来说服你：这个概念并不怎么陌生。

\trSheaf{}之所以滑不留手，原因之一在于它们记录着海量的信息，其间还有一些微妙的局部到整体的
问题。此外，把握\trSheaf{}共有三条不同的途径：
\begin{itemize}
\item 以\trOpenSet{}表述（即 \S2.2 的定义）——直观，但在某些方面最不管用；
\item 以\trStalk{}表述（见 \S2.4.1）；以及
\item 以拓扑的\trBase{}表述（\S2.5）。
\end{itemize}
（也有些人强烈偏爱\trEspaceEtale{}（espace étalé）的解释，见 \S2.2.11。）何者当用须凭经验，
因此务必做相当数量的、涉及\trSheaf{}的不同侧面的练习，才能真正理解这个概念。
